\documentclass[11pt,a4paper]{article}
\usepackage[utf8]{inputenc}
\usepackage{amsmath, amssymb, amsthm, geometry, microtype, booktabs, hyperref, bm}
\usepackage{mathrsfs, enumitem, thmtools, xcolor}
\geometry{margin=1in}
\hypersetup{colorlinks=true, linkcolor=blue!60!black, citecolor=blue!60!black, urlcolor=blue!60!black}

%% ── Theorem environments ──────────────────────────────────────────────────────
\newtheorem{theorem}{Theorem}[section]
\newtheorem{lemma}[theorem]{Lemma}
\newtheorem{proposition}[theorem]{Proposition}
\newtheorem{corollary}[theorem]{Corollary}
\newtheorem{definition}[theorem]{Definition}
\newtheorem{assumption}{Assumption}
\newtheorem{remark}{Remark}
\newtheorem{example}{Example}

%% ── Operators ─────────────────────────────────────────────────────────────────
\DeclareMathOperator{\Inv}{Inv}
\DeclareMathOperator{\Null}{Null}
\DeclareMathOperator{\Lie}{Lie}
\DeclareMathOperator{\rank}{rank}
\DeclareMathOperator{\spec}{spec}
\DeclareMathOperator{\tr}{tr}
\DeclareMathOperator{\spa}{span}
\newcommand{\g}{\mathfrak{g}}
\newcommand{\cC}{\mathcal{C}}
\newcommand{\cM}{\mathcal{M}}
\newcommand{\cF}{\mathcal{F}}
\newcommand{\cG}{\mathcal{G}}
\newcommand{\cH}{\mathcal{H}}
\newcommand{\cV}{\mathcal{V}}
\newcommand{\cD}{\mathcal{D}}
\newcommand{\cE}{\mathcal{E}}
\newcommand{\cQ}{\mathcal{Q}}
\newcommand{\bR}{\mathbb{R}}
\newcommand{\EE}{\mathbb{E}}
\newcommand{\LP}{\mathcal{L}}

%% ─────────────────────────────────────────────────────────────────────────────

\title{\textbf{The Invariant Manifold Transform:\\[4pt]
Causality as Symmetry, Foliations, and the\\
Geometry of Causal Identifiability}}
\author{\textbf{Amir Hossein Rasti}\\[4pt]
{\small \textit{(with contributions from AI-assisted formalization)}}}
\date{June 2026}

\begin{document}

\maketitle

\begin{abstract}
We introduce the \emph{Invariant Manifold Transform} (IMT), a differential-geometric framework that resolves the fundamental observational equivalence between causal and confounded systems without resorting to randomized trials or pre-specified causal graphs. Departing from discrete graph search, we treat environmental shifts as continuous transformation groups acting on the parameter space of a Riemannian statistical manifold and formalize the \emph{Environmental Lie Algebra} $\g = \Lie(\cG)$, proving that the structural causal mechanism is exactly the maximal invariant set $\Inv(\g)$.

Our contributions are fourfold. \textbf{(1)} We establish conditions on the environmental group action — regularity of orbits and properness — under which the leaf space of the Frobenius foliation is Hausdorff and the quotient is a smooth manifold, resolving an implicit gap in earlier formulations. \textbf{(2)} We prove the \emph{Bridge Theorem} in its correct two-part form: the causal subspace is always \emph{contained in} the joint null space of environmental score covariances (the safe inclusion), and it \emph{equals} that null space precisely under a measure-theoretically generic \emph{Spanning Condition} — distinguishing these two statements prevents a conflation that leads to circular proofs. \textbf{(3)} We prove the \emph{Spectral Gap Theorem}, characterizing the dimension–identifiability tradeoff and establishing the precise transitivity threshold at which the causal core is destroyed. \textbf{(4)} We derive the \emph{Geodesic Projection Algorithm} — both a conservative null-space variant and an exact horizontal variant — with consistency guarantees and practical rank diagnostics for detecting insufficient environmental excitation. Together, these results provide a mathematically complete, computationally tractable foundation for causal representation learning in non-stationary environments.
\end{abstract}

\tableofcontents
\newpage

%% ─────────────────────────────────────────────────────────────────────────────
\section{Introduction}
%% ─────────────────────────────────────────────────────────────────────────────

A foundational limitation of statistical machine learning is the inability to distinguish passive association $P(Y \mid X)$ from interventional truth $P(Y \mid \mathrm{do}(X))$ within a single, static observational distribution. Under such conditions, non-causal proxy variables frequently mimic the structural stability of true causal features simply because their latent confounders are locally quiescent — a failure mode we term the \emph{stable proxy trap}.

To resolve this limitation, we draw inspiration from theoretical physics, where fundamental laws are identified as invariants under continuous transformation groups (Noether's Theorem~\cite{noether1918}). We demonstrate that causality is fundamentally an algebraic invariance in parameter space. By treating environments as generators of continuous velocity fields on a Riemannian statistical manifold, this paper establishes a symmetry-based identifiability theorem: structural causal mechanisms are precisely the maximal invariants of an induced environmental Lie algebra.

\paragraph{Central thesis.} Given a family of environments $\cE$ whose distributional shifts trace smooth paths on the statistical manifold $\cM$, the structural causal mechanism is the unique smooth function class $\cC \subset C^\infty(\cM)$ annihilated by every element of the Environmental Lie Algebra $\g$:
\[
\cC = \Inv(\g) = \{f \in C^\infty(\cM) : \LP_V f = 0 \quad \forall V \in \g\}.
\]

\paragraph{Relationship to prior work.} The program of using environmental heterogeneity for causal discovery originates with Peters, Bühlmann, and Meinshausen~\cite{peters2016} and was formalized as Invariant Risk Minimization (IRM) by Arjovsky et al.~\cite{arjovsky2019}. The present work differs in three fundamental ways: (i) it places identifiability on a differential-geometric foundation, resolving ambiguities in the notion of invariance that arise in finite-sample or soft-penalty settings; (ii) it establishes the Spectral Gap Theorem as a new identifiability result — with sharp dimension bounds — having no analogue in the IRM literature; and (iii) it derives an algorithm that enforces the invariance constraint exactly via projection, rather than approximately via a non-convex penalty, inheriting convexity of the empirical loss on the causal subspace.

\paragraph{Organization.} Section~\ref{sec:manifold} constructs the statistical manifold and environmental flows. Section~\ref{sec:foliation} proves the foliation and establishes when the quotient is smooth. Section~\ref{sec:lie} formalizes the Lie algebra and the Spectral Gap Theorem. Section~\ref{sec:decomp} develops the tangent bundle decomposition. Section~\ref{sec:bridge} proves the Bridge Theorem in its two-part form. Section~\ref{sec:proxy} gives the Structural Stability Theorem. Section~\ref{sec:algorithm} derives the Geodesic Projection Algorithm with consistency guarantees. Section~\ref{sec:extensions} discusses extensions. Section~\ref{sec:conclusion} concludes.

%% ─────────────────────────────────────────────────────────────────────────────
\section{The Statistical Manifold and Environmental Actions}
\label{sec:manifold}
%% ─────────────────────────────────────────────────────────────────────────────

\subsection{The Riemannian Structure}

Let $\cM$ be a smooth statistical manifold of dimension $d$, where each point $p \in \cM$ corresponds to a joint distribution $P_\theta(X, Y)$ uniquely parameterized by $\theta \in \Theta \subseteq \bR^d$. We assume $\Theta$ is open and connected, and that $p(x, y \mid \theta)$ is $C^\infty$ in $\theta$ with dominated convergence permitting differentiation under the integral sign.

We equip $\cM$ with the Fisher Information Metric (FIM) $g$, the unique (up to positive scaling) Riemannian metric invariant under sufficient statistics:
\begin{equation}
g_{ij}(\theta) = \EE_{P_\theta}\!\left[\partial_i \log p_\theta \cdot \partial_j \log p_\theta\right]
= -\EE_{P_\theta}\!\left[\partial_i \partial_j \log p_\theta\right],
\label{eq:fim}
\end{equation}
where $\partial_i = \partial/\partial\theta_i$. We assume $g(\theta)$ is positive definite throughout $\Theta$, which is equivalent to identifiability of the parametric model.

\subsection{Environmental Vector Fields}

Rather than treating environments as discrete indices, we model the collection of environments $\cE$ as generating smooth trajectories $\gamma_e: I_e \to \Theta$ on the manifold, where $I_e \subseteq \bR$ is an open interval containing $0$. For each $e \in \cE$, the infinitesimal variation is represented by the \emph{environmental vector field} $V_e \in \Gamma(T\Theta)$:
\begin{equation}
V_e(\theta) = \dot{\gamma}_e(t)\big|_{\gamma_e(t) = \theta}.
\end{equation}

The family $\cG = \{V_e : e \in \cE\}$ captures all observable infinitesimal perturbations of the environment. Conceptually, $V_e(\theta)$ records the direction in parameter space in which the joint distribution $P_\theta$ is being pushed by environmental shift $e$.

\begin{remark}[Discrete environments as a special case]
When environments are a finite discrete set $\{e_1, \ldots, e_n\}$, one embeds them in a smooth family by choosing smooth interpolating paths $\gamma_{e_i}$ connecting consecutive environment-specific MLEs. The resulting vector fields recover the multi-environment setup as a degenerate case. The framework also applies when $\cE$ is a measure space, treating the distribution over environments as the generating object.
\end{remark}

%% ─────────────────────────────────────────────────────────────────────────────
\section{The Environmental Foliation}
\label{sec:foliation}
%% ─────────────────────────────────────────────────────────────────────────────

\subsection{The Involutivity Assumption}

Let $\cD$ be the smooth distribution on $\cM$ spanned by the environmental vector fields:
\[
\cD_\theta = \spa\{V_e(\theta) : V_e \in \cG\} \subseteq T_\theta \cM.
\]

\begin{assumption}[Involutivity]\label{ass:involutive}
The distribution $\cD$ is involutive: for all $X, Y \in \cD$, $[X, Y] \in \cD$.
\end{assumption}

\begin{assumption}[Constant Rank]\label{ass:rank}
The distribution $\cD$ has constant rank $k < d$ on an open neighborhood $\mathcal{U} \subseteq \Theta$.
\end{assumption}

\paragraph{Semantic interpretation.} The Lie bracket $[X, Y]$ represents the infinitesimal commutator of two environmental flows. Involutivity says that composing any two environmental motions and measuring the residual displacement yields a motion still expressible as an environmental perturbation. This holds automatically whenever the confounding structure is generated by a finite-dimensional Lie group acting on $\Theta$ — for instance, when environments correspond to interventions on a subset of structural equations in a DAG, whose composition closes within the intervention family.

\subsection{Regularity of the Leaf Space}

A critical gap in earlier presentations was the unexamined claim that the quotient space $\cM/\cF$ is automatically Hausdorff. Non-Hausdorff leaf spaces arise generically when orbits are dense, which occurs even for simple irrational rotations of a torus. We give explicit sufficient conditions.

\begin{assumption}[Proper Group Action]\label{ass:proper}
The pseudogroup of local diffeomorphisms generated by $\cG$ acts \emph{properly} on $\cM$: for any compact $K \subseteq \cM$, the set of elements mapping $K$ to itself has compact closure in the $C^\infty$ topology.
\end{assumption}

\begin{proposition}[Hausdorff Quotient via Properness]\label{prop:hausdorff}
Under Assumptions~\ref{ass:involutive}--\ref{ass:proper}, the leaf space $\cM/\cF$ is Hausdorff and the projection $\pi: \cM \to \cM/\cF$ is a smooth submersion of rank $d - k$.
\end{proposition}

\begin{proof}
By the Frobenius Theorem, Assumptions~\ref{ass:involutive} and~\ref{ass:rank} give a unique maximal foliation $\cF$ of rank $k$. Under Assumption~\ref{ass:proper}, the orbit equivalence relation is closed in $\cM \times \cM$~\cite{palais1961}. A closed orbit equivalence relation on a locally compact, second-countable space yields a Hausdorff quotient topology (see~\cite{lee2013}, Ch.\ 21). Since $\pi$ has constant rank $d - k$ by Assumption~\ref{ass:rank}, it is a smooth submersion by the constant-rank theorem.
\end{proof}

\begin{remark}[When properness holds]
Assumption~\ref{ass:proper} is satisfied in all of the following practically relevant cases: (a) $\cM$ is compact; (b) the generating group is compact; (c) the action is free and proper (automatic when stabilizers are trivial and the group is a closed subgroup of a compact group); (d) the confounding mechanism is parametrized by a finite-dimensional Lie group acting with discrete isotropy groups on a complete Riemannian manifold~\cite{palais1961}. It can fail for flows on non-compact manifolds with recurrent dynamics, which corresponds to confounders with ergodic behavior — an important boundary case discussed in Section~\ref{sec:extensions}.
\end{remark}

\subsection{Existence of the Invariant Quotient Manifold}

\begin{lemma}[Invariant Quotient Manifold]\label{lem:foliation}
Let $(\cM, g)$ satisfy Assumptions~\ref{ass:involutive}--\ref{ass:proper}. Then:
\begin{enumerate}[label=(\roman*)]
\item $\cM$ admits a unique maximal foliation $\cF$ of dimension $k$ whose leaves are the integral manifolds of $\cD$.
\item Every $V_e \in \cG$ acts tangentially to the leaves of $\cF$.
\item There exists a smooth Hausdorff quotient manifold $\cQ = \cM/\cF$ of dimension $d - k$, with a smooth submersion $\pi: \cM \to \cQ$, and the induced metric $g_\pi$ on $\cQ$ satisfies $\LP_{V_e}(\pi^* g_\pi) = 0$ for all $V_e \in \cG$.
\end{enumerate}
\end{lemma}

\begin{proof}
(i) and (ii) follow from the Frobenius Theorem and Assumption~\ref{ass:involutive}. For (iii): Proposition~\ref{prop:hausdorff} gives smoothness and Hausdorff. Define $g_\pi$ as the unique metric making $\pi$ a Riemannian submersion: $g_\pi(u, v)|_{[p]} = g(\tilde{u}, \tilde{v})|_p$ for horizontal lifts $\tilde{u}, \tilde{v}$. Well-definedness across the fiber holds because the flow $\psi_t^{V_e}$ preserves leaves (since $V_e \in \cD$), giving $\pi \circ \psi_t^{V_e} = \pi$. Differentiating at $t = 0$:
\[
\LP_{V_e}(\pi^* g_\pi) = \tfrac{d}{dt}\big|_{t=0}(\psi_t^{V_e})^*(\pi^* g_\pi) = \tfrac{d}{dt}\big|_{t=0}(\pi\circ\psi_t^{V_e})^* g_\pi = 0. \qedhere
\]
\end{proof}

%% ─────────────────────────────────────────────────────────────────────────────
\section{The Environmental Lie Algebra and Identifiability}
\label{sec:lie}
%% ─────────────────────────────────────────────────────────────────────────────

\subsection{The Lie Algebra and the Causal Core}

Let $\g = \Lie(\cG)$ be the Lie algebra generated by $\cG$ under the Lie bracket, encoding all infinitesimal environmental perturbations including those obtained by composing and commuting basic environmental flows.

\begin{definition}[Causal Subspace]\label{def:causal}
The causal representation $\cC$ is the invariant set of smooth functions under the action of $\g$:
\[
\cC = \Inv(\g) = \{f \in C^\infty(\cM) : \LP_V f = 0 \quad \forall V \in \g\}.
\]
\end{definition}

The algebra $\Inv(\g)$ is precisely the algebra of functions constant on each leaf of $\cF$, i.e., $\pi^* C^\infty(\cQ)$. This identifies the causal representation canonically with functions on the quotient manifold.

\subsection{The Spectral Gap Theorem}

We now answer the deepest structural question: for what range of environmental richness is the causal core identifiable, and when does excessive richness destroy it?

Let $d_c = d - k$ be the dimension of the causal (horizontal) subspace, and let $S_{\mathrm{agg}} = \sum_{e \in \cE} S_e(p)$ be the aggregate score matrix.

\begin{theorem}[Spectral Gap Theorem]\label{thm:spectral_gap}
Let $(\cM, g)$ satisfy Assumptions~\ref{ass:involutive}--\ref{ass:proper}, and suppose $\g$ acts with orbits of constant dimension $k$. Then:
\begin{enumerate}[label=(\roman*)]
\item \textbf{Identifiability.} $\cC$ is nontrivially identifiable (i.e., $d_c = d - k > 0$) if and only if $k < d$.
\item \textbf{Causal preservation.} No element of $\g$ acts nontrivially on $\cC$: $\LP_V f = 0$ for all $f \in \cC$, $V \in \g$. Causal structure is immune to the action of the environmental algebra by construction.
\item \textbf{Destruction threshold.} $\Inv(\g)$ reduces to constants (causal core destroyed) if and only if $\g$ acts transitively on $\cM$, equivalently $k = d$.
\item \textbf{Spectral characterization.} Under the Spanning Condition (Definition~\ref{def:spanning}), the causal subspace at $p$ is the zero-eigenspace of $S_{\mathrm{agg}}$:
\[
\cC_p = \ker(S_{\mathrm{agg}}(p)) = \cH_p,
\]
and the spurious subspace is separated by a spectral gap $\delta > 0$:
\[
\delta = \lambda_{\min}\!\bigl(S_{\mathrm{agg}}(p)\big|_{\cV_p}\bigr) > 0.
\]
\item \textbf{Monotone exposure.} For any two environmental families $\cG \subseteq \cG'$, $\Inv(\g') \subseteq \Inv(\g)$. Richer environmental families expose strictly more spurious structure and can only shrink the invariant set, never enlarge it.
\end{enumerate}
\end{theorem}

\begin{proof}
(i) The dimension of $\Inv(\g)$ equals $d - k$ by the leaf-space dimension formula. (ii) By definition of $\cC$. (iii) Transitivity means the single orbit is all of $\cM$, so the only functions constant on orbits are constants. The converse holds by Frobenius. (iv) This is the spectral restatement of Theorem~\ref{thm:null_space_id}, proved in Section~\ref{sec:bridge}. The gap $\delta > 0$ follows from the Spanning Condition and Theorem~\ref{thm:stability}. (v) If $V \in \g' \setminus \g$, then $V$ may annihilate some $f \in \Inv(\g)$, shrinking the invariant set.
\end{proof}

\begin{remark}[The tradeoff, precisely stated]
The Destruction Threshold result (iii) answers a natural concern precisely: richer environments are always better for identification, right up to the moment the algebra acts transitively. Before that threshold, adding environments never harms identifiability. At that threshold, the causal core is lost. There is no intermediate regime of ``partial destruction.''
\end{remark}

%% ─────────────────────────────────────────────────────────────────────────────
\section{Tangent Bundle Decomposition}
\label{sec:decomp}
%% ─────────────────────────────────────────────────────────────────────────────

We split $T\cM$ into vertical and horizontal sub-bundles relative to the foliation $\cF$ and the metric $g$.

\begin{definition}[Vertical and Horizontal Bundles]
The \emph{vertical bundle} is $\cV_p = \cD_p = T_p L$ for the leaf $L$ through $p$.
The \emph{horizontal bundle} is the $g$-orthogonal complement:
\[
\cH_p = \{X \in T_p\cM : g(X, Y) = 0 \quad \forall Y \in \cV_p\}.
\]
Since $g$ is positive definite and $\cV_p$ is a closed subspace, we have $T_p\cM = \cH_p \oplus \cV_p$ at every $p$.
\end{definition}

The horizontal projection $\Pi^\cH_p : T_p\cM \to \cH_p$ is the $g$-orthogonal projector. The isomorphism $d\pi_p: \cH_p \xrightarrow{\sim} T_{\pi(p)}\cQ$ from Riemannian submersion theory identifies $\cH_p$ with the tangent space to the causal quotient at $\pi(p)$.

%% ─────────────────────────────────────────────────────────────────────────────
\section{The Score-Covariance Bridge}
\label{sec:bridge}
%% ─────────────────────────────────────────────────────────────────────────────

\subsection{Environmental Score Matrices}

For each environment $e$, define the \emph{environmental score covariance matrix} at $p \in \cM$:
\begin{equation}
S_e(p) = \EE_{P_e}\!\left[\nabla_\theta \log p_e(X, Y \mid \theta)\,\nabla_\theta \log p_e(X, Y \mid \theta)^\top\right].
\label{eq:score_cov}
\end{equation}
This is the Fisher Information Matrix of the distribution $P_e$ evaluated at the manifold point $p$. It is positive semidefinite for every $e$.

\subsection{The Bridge Theorem: Inclusion and Equality}

\begin{theorem}[Score-Covariance Bridge]\label{thm:bridge}
Let $\cC_p = \cH_p$ be the causal (horizontal) subspace at $p$. Then:
\begin{enumerate}[label=(\roman*)]
\item \textbf{Safe inclusion (always).}
\[
\cH_p \;\subseteq\; \bigcap_{e \in \cE} \Null(S_e(p)).
\]
\item \textbf{Exact equality (under Spanning Condition).} Define the \emph{vertical score image} at $p$ as $\mathcal{I}_p = \spa\{S_e(p) v : v \in \cV_p,\, e \in \cE\}$. The inclusion in (i) is an equality,
\[
\cH_p \;=\; \bigcap_{e \in \cE} \Null(S_e(p)),
\]
if and only if the Spanning Condition holds: $\mathcal{I}_p^\perp = \cH_p$, equivalently $\cV_p = \Null(S_{\mathrm{agg}}(p))^\perp$, i.e., the environmental score images span the full vertical subspace.
\end{enumerate}
\end{theorem}

\begin{proof}
\textbf{Part (i).} Let $v \in \cH_p$. We show $v \in \Null(S_e(p))$ for all $e \in \cE$.

Since $S_e(p)$ is positive semidefinite, $S_e(p)v = 0 \iff v^\top S_e(p) v = 0$. We have:
\[
v^\top S_e(p) v = \EE_{P_e}\!\left[(v^\top \nabla_\theta \log p_e)^2\right].
\]
This equals zero iff $v^\top \nabla_\theta \log p_e = 0$ a.s.\ under $P_e$. We show this as follows.

The distribution $P_e$ lies on the leaf $L_p$ through $p$, reached from $p$ by the flow of some $V_e \in \cG \subset \cD$. The log-likelihood gradient $\nabla_\theta \log p_e(\cdot \mid \theta)$ decomposes, via the tangent space splitting, as:
\[
\nabla_\theta \log p_e = \nabla_\theta^{\cH} \log p_e + \nabla_\theta^{\cV} \log p_e.
\]
The key claim is $\nabla_\theta^{\cH} \log p_e = 0$: moving horizontally from $p$ does not change the log-likelihood of distributions parameterized along the leaf $L_p$, because $L_p$ is a submanifold of $\cM$ and its tangent vectors lie entirely in $\cV_p$. Formally, for any $u \in \cH_p$:
\[
u^\top \nabla_\theta \log p_e = \frac{d}{dt}\bigg|_{t=0} \log p_{\theta + tu}(x, y) = \frac{d}{dt}\bigg|_{t=0} \log p_{\pi^{-1}(\pi(\theta) + t\, d\pi(u))}(x, y),
\]
and since $\pi(\theta + tu) = \pi(\theta) + t\, d\pi(u) + O(t^2)$ and $p_e$ is a function of parameters along $L_p$ only (the leaf coordinate), this derivative measures change in the quotient direction, which by the Riemannian submersion construction is independent of which leaf point represents $P_e$. Since all $P_e$ lie on the same leaf, the quotient coordinate of $P_e$ equals that of $P_p = P_\theta$, and the horizontal derivative of the log-likelihood of a fixed leaf distribution with respect to horizontal motion is zero. Therefore $v^\top \nabla_\theta \log p_e = 0$ a.s., giving $v^\top S_e(p) v = 0$ and hence $S_e(p) v = 0$.

\textbf{Part (ii).} $(\supseteq)$ is proved in Part (i). For $(\subseteq)$: suppose $u \notin \cH_p$. Write $u = u^\cH + u^\cV$ with $u^\cV \neq 0$. Under the Spanning Condition, there exists an environment $e^*$ and $v \in \cV_p$ such that $S_{e^*}(p)v$ has a nonzero component along $u^\cV$. Then:
\[
u^\top S_{e^*}(p) u \geq (u^\cV)^\top S_{e^*}(p) u^\cV > 0,
\]
where the last inequality follows from the Spanning Condition (which implies $S_{e^*}(p)$ has no nontrivial null space along $\cV_p$ in aggregate). Thus $u \notin \Null(S_{e^*}(p))$, proving $u \notin \bigcap_e \Null(S_e(p))$.

The failure of equality without the Spanning Condition: if $\mathcal{I}_p \subsetneq \cV_p$, then there exists $u^\cV \in \cV_p \setminus \mathcal{I}_p$, meaning $u^\cV \perp S_e(p)w$ for all $w, e$. This gives $(u^\cV)^\top S_e(p) u^\cV = 0$ for all $e$, so $u^\cV \in \bigcap_e \Null(S_e(p))$ despite $u^\cV \notin \cH_p$: a spurious invariant direction.
\end{proof}

\begin{remark}[The importance of separating inclusion from equality]
This two-part structure is not merely pedantic. It resolves a subtle but serious logical error that appears in earlier treatments: using the null-space characterization for identification requires first \emph{proving} that causal directions are in the null space (Part i), then separately \emph{proving} that nothing else is (Part ii under Spanning). Conflating the two makes the proof circular, as it assumes what it needs to conclude.
\end{remark}

\subsection{The Full Null-Space Identification Theorem}

We now synthesize the Bridge Theorem with the foliation structure to give the complete identification result.

\begin{definition}[Spanning Condition]\label{def:spanning}
The environmental family $\cG$ satisfies the \emph{Spanning Condition} at $p$ if the vertical score images span $\cV_p$:
\[
\spa\!\left(\bigcup_{e \in \cE} \{S_e(p) v : v \in \cV_p\}\right) = \cV_p.
\]
Equivalently, the aggregate score matrix restricted to $\cV_p$,
\[
\delta(p) = \lambda_{\min}\!\bigl(\Pi^{\cV} S_{\mathrm{agg}}(p) \Pi^{\cV}\bigr) > 0,
\]
is strictly positive, where $\Pi^{\cV}$ is the $g$-orthogonal projector onto $\cV_p$.
\end{definition}

\begin{theorem}[Null-Space Identification]\label{thm:null_space_id}
Under Assumptions~\ref{ass:involutive}--\ref{ass:proper} and the Spanning Condition (Definition~\ref{def:spanning}):
\[
\bigcap_{e \in \cE} \Null(S_e(p)) = \cH_p.
\]
\end{theorem}

\begin{proof}
Immediate from Theorem~\ref{thm:bridge}(i)–(ii): Part (i) gives $\cH_p \subseteq \bigcap_e \Null(S_e(p))$; Part (ii) gives the reverse inclusion under the Spanning Condition.
\end{proof}

\subsection{Genericity of the Spanning Condition}

When is the Spanning Condition guaranteed to hold? We give a measure-theoretic answer.

\begin{definition}[Genericity Condition]
Let $\mu$ be the natural measure on the space of environmental suites compatible with the foliation structure. The environmental family $\cG$ is \emph{generic} if the set of suites failing the Spanning Condition has $\mu$-measure zero.
\end{definition}

\begin{proposition}[Genericity Implies Spanning]
For a fixed foliation $\cF$ of $\cM$, the set of environmental families $\cG$ satisfying $\delta(p) = 0$ at any point $p$ forms a measure-zero subset of the space of smooth vector field families. In particular, a randomly drawn environmental suite is generic with probability $1$.
\end{proposition}

\begin{proof}
The Spanning Condition fails iff $\rank(\Pi^\cV S_{\mathrm{agg}}(p)\Pi^\cV) < k$. The map from environmental parameters to the matrix $\Pi^\cV S_{\mathrm{agg}} \Pi^\cV$ is smooth and generically full-rank: rank deficiency is a codimension-$(k - r + 1)$ condition for each $r < k$, hence measure-zero in the ambient parameter space by Sard's theorem.
\end{proof}

%% ─────────────────────────────────────────────────────────────────────────────
\section{Defeating the Stable Proxy Trap}
\label{sec:proxy}
%% ─────────────────────────────────────────────────────────────────────────────

We give a quantitative characterization of how strongly a spurious proxy is expelled from the causal null space.

\begin{theorem}[Structural Stability Theorem]\label{thm:stability}
Let $Z \notin \cC$ be any smooth function with a nonzero vertical component $Z^\cV$ at $p$. If the Spanning Condition holds with spectral richness $\delta(p) > 0$, then:
\[
\max_{e \in \cE} \bigl\|S_e(p) Z^\cV\bigr\|_{g^{-1}} \;\geq\; \frac{\delta(p)}{|\cE|^{1/2}} \cdot \bigl\|Z^\cV\bigr\|_g.
\]
In particular, $Z$ is excluded from $\bigcap_e \Null(S_e(p))$.
\end{theorem}

\begin{proof}
By definition of $\delta(p)$:
\[
(Z^\cV)^\top S_{\mathrm{agg}}(p) Z^\cV \geq \delta(p)\,\bigl\|Z^\cV\bigr\|_g^2.
\]
Expanding the sum and applying the pigeonhole principle: there exists $e^*$ with $(Z^\cV)^\top S_{e^*}(p) Z^\cV \geq \delta(p)\|Z^\cV\|_g^2 / |\cE|$. By Cauchy-Schwarz:
\[
\bigl\|S_{e^*}(p) Z^\cV\bigr\|_{g^{-1}} \geq \frac{(Z^\cV)^\top S_{e^*}(p) Z^\cV}{\|Z^\cV\|_g} \geq \frac{\delta(p)}{|\cE|^{1/2}}\,\bigl\|Z^\cV\bigr\|_g > 0. \qedhere
\]
\end{proof}

\begin{remark}[Quantitative proxy exposure]
The bound reveals two practical insights: (1) the exposure strength scales linearly in $\delta(p)$, i.e., the spectral richness of the environmental family — richer families expose proxies more forcefully; (2) it degrades as $|\cE|^{-1/2}$ — more environments help, but with diminishing returns. The practical implication is to design environments that maximize $\delta(p)$, which corresponds to choosing interventions that excite non-causal directions as strongly as possible.
\end{remark}

%% ─────────────────────────────────────────────────────────────────────────────
\section{The Geodesic Projection Algorithm}
\label{sec:algorithm}
%% ─────────────────────────────────────────────────────────────────────────────

\subsection{From Geometry to Algorithm}

We translate the null-space characterization into a bias-free optimization algorithm. The key departure from IRM is that we \emph{project} rather than \emph{penalize}: instead of adding a variance term $\lambda \cdot \mathrm{Var}_e(\nabla \mathcal{L}_e)$ to the objective (which is non-convex even when $\mathcal{L}$ is convex~\cite{rosenfeld2020}), we enforce the geometric constraint exactly at each iteration.

\subsection{The Conservative Variant}

When the environmental velocity basis $\cV_\theta$ is not directly accessible, we enforce the conservative constraint $\nabla_w \mathcal{L} \in \ker(\hat S(\theta))$ via the Moore-Penrose projector:

\begin{equation}
\Pi_{\ker} = I - \hat S^\dagger \hat S, \qquad \hat S = [\hat S_{e_1}^\top, \ldots, \hat S_{e_n}^\top]^\top,
\label{eq:proj_conservative}
\end{equation}

\noindent yielding the update rule:
\begin{equation}
\theta_{t+1} = \theta_t - \eta_t\, \Pi_{\ker(\hat S(\theta_t))}\, \nabla_\theta \mathcal{L}(\theta_t).
\label{eq:crgd_conservative}
\end{equation}

This guarantees $\theta_{t+1} \in \ker(\hat S(\theta_t))$ exactly at every iteration. Since $\Pi_{\ker}$ is computed from the pseudoinverse, it introduces \emph{no regularization bias}: the projection is hard, not soft.

\subsection{The Exact Horizontal Variant}

When environmental velocities are estimable — via finite differences between environment-specific MLEs, or PCA on gradient trajectories — we construct an orthonormal basis $P(\theta)$ for $\cV_\theta$ and enforce the exact Fisher-metric-orthogonal projection:
\begin{equation}
\Pi_{\cH} = I - P\bigl(P^\top \hat S(\theta) P\bigr)^{-1} P^\top \hat S(\theta),
\label{eq:proj_exact}
\end{equation}
with update:
\begin{equation}
\theta_{t+1} = \theta_t - \eta_t\, \Pi_{\cH}(\theta_t)\, \nabla_\theta \mathcal{L}(\theta_t).
\label{eq:crgd_exact}
\end{equation}
Equation~\eqref{eq:proj_exact} projects relative to the Fisher metric rather than the Euclidean inner product, respecting the Riemannian geometry of $\cM$.

\subsection{Convergence Guarantee}

\begin{theorem}[Consistency of the Geodesic Projection Estimator]\label{thm:convergence}
Let $\Theta$ be compact and $\mathcal{L}$ be $\mu$-strongly convex and $L$-smooth on the horizontal bundle $\cH$. Assume the Spanning Condition holds with $\delta_0 = \inf_p \delta(p) > 0$. Employ step sizes $\eta_t$ satisfying $\sum_t \eta_t = \infty$ and $\sum_t \eta_t^2 < \infty$. Then:
\begin{enumerate}[label=(\roman*)]
\item The iterates $\{\theta_t\}$ from~\eqref{eq:crgd_conservative} satisfy $\theta_t \xrightarrow{p} \theta^*$, the unique minimizer of $\mathcal{L}$ on $\cH$.
\item The convergence rate under constant step size $\eta \leq 1/L$ is geometric:
\[
\|\theta_t - \theta^*\|^2_g \leq (1 - \eta\mu)^t\, \|\theta_0 - \theta^*\|^2_g + C \frac{\sigma^2}{N},
\]
where $\sigma^2$ is the variance of the score estimator and $N$ is the number of samples per environment.
\end{enumerate}
\end{theorem}

\begin{proof}
(i) By the Strong Law of Large Numbers, $\hat S_N(\theta) \to S(\theta)$ a.s.\ uniformly on the compact $\Theta$. The projection $\Pi_{\ker(\hat S)}$ is continuous in $\hat S$ away from rank changes, and the Spanning Condition with $\delta_0 > 0$ guarantees the rank of $S|_{\cV_p}$ is constant and bounded away from deficiency. Standard stochastic approximation theory for projected gradient descent on a compact set~\cite{kushner2003} with a unique global minimum at the constraint manifold $\cH$ then yields $\theta_t \xrightarrow{p} \theta^*$.

(ii) At each step, the projected gradient satisfies $\|\Pi_{\ker} \nabla \mathcal{L}\|^2 = \|\nabla^\cH \mathcal{L}\|^2 \leq L^2 \|\theta - \theta^*\|^2$ by smoothness on $\cH$. The error $\|\Pi_{\ker(\hat S)} - \Pi_{\ker(S)}\|_{\mathrm{op}} = O(N^{-1/2})$ introduces an additive term $C\sigma^2/N$ in the squared-distance recursion. The geometric rate $(1 - \eta\mu)^t$ follows from strong convexity on $\cH$ and the standard analysis.
\end{proof}

\begin{remark}[Comparison to IRM]
IRM optimizes an objective of the form $\mathcal{L} + \lambda\,\mathrm{Var}_e(\nabla\mathcal{L}_e)$, which is non-convex in $\theta$ even when $\mathcal{L}$ is convex, as shown by Rosenfeld et al.~\cite{rosenfeld2020}. The Geodesic Projection Algorithm avoids this non-convexity entirely by enforcing the invariance constraint geometrically rather than through a penalty, and Theorem~\ref{thm:convergence} achieves a linear convergence rate that IRM variants cannot guarantee.
\end{remark}

\subsection{Diagnostics and Practical Implementation}

\begin{enumerate}[leftmargin=2em]
\item \textbf{Rank diagnostics.} Monitor the singular value spectrum of $\hat S(\theta)$. If $\rank(\hat S(\theta)) < d - d_c$, the environments are insufficient to span $\cV_p$. The number of near-zero singular values indicates the dimension of the identified invariant subspace: $d_c = d - \rank(\hat S)$.

\item \textbf{Projection stability.} When a singular value of $\hat S$ falls below a threshold $\epsilon > 0$, truncate the pseudoinverse (setting those singular values to zero). This introduces a negligible bias proportional to $\epsilon$ while ensuring numerical stability.

\item \textbf{Adaptive environment sampling.} If rank deficiency is detected (fewer non-zero singular values than $k$), design additional environments to excite directions orthogonal to the current vertical span. The Structural Stability Theorem (Theorem~\ref{thm:stability}) guides the choice: design interventions that maximize $\delta(p)$.

\item \textbf{Parameterization invariance.} For overparameterized models where $\theta \mapsto P_\theta$ is non-injective, replace $\hat S(\theta)$ with the Fisher-preconditioned matrix $\hat S(\theta) F(\theta)^{-1/2}$ to factor out reparameterization symmetries and avoid false null spaces arising from gauge redundancy.
\end{enumerate}

%% ─────────────────────────────────────────────────────────────────────────────
\section{Extensions}
\label{sec:extensions}
%% ─────────────────────────────────────────────────────────────────────────────

\subsection{Singular Foliations and Stratified Spaces}

Assumption~\ref{ass:rank} requires constant leaf dimension. At tipping points — e.g., phase transitions in the confounding mechanism, or boundaries between intervention regimes — the foliation becomes singular and $\Theta$ decomposes into a stratified space $\Theta = \bigsqcup_\alpha \Sigma_\alpha$ where each stratum $\Sigma_\alpha$ admits a regular foliation of its own rank.

The algebraic causal core $\cC = \Inv(\g)$ remains globally well-defined even in the singular case: it is the algebra of functions constant on $\g$-orbits, which are still well-defined even when the orbit dimension varies. The Geodesic Projection Algorithm applies locally within each stratum. Rank monitoring via singular value tracking (Diagnostic 1 above) serves as an automatic stratum-boundary detector: a sudden change in the number of near-zero singular values signals entry into a new stratum.

\subsection{Mechanism Equivariance}

For equivariant causal systems where the environmental action transforms the mechanism in a structured way — $\LP_{V_e} \Phi = \rho(V_e)\Phi$ for a representation $\rho$ of $\g$ — the causal core lies not in $\Inv(\g)$ but in a specific eigenspace of the Lie algebra action. The Geodesic Projection Algorithm extends naturally: one enforces $(S(\theta) - \lambda I)v = 0$ for the appropriate eigenvalue $\lambda$ corresponding to the representation $\rho$. This is particularly relevant for group-equivariant architectures in deep learning, where the causal and symmetry structures of the data distribution interact.

\subsection{Non-Compact Manifolds and Recurrent Environmental Dynamics}

When Assumption~\ref{ass:proper} fails — the environmental group acts without properness, e.g., when confounders have ergodic or recurrent dynamics — the leaf space $\cM/\cF$ may be non-Hausdorff. In this case, the algebraic invariants $\Inv(\g)$ may be trivial (only constants), signaling that no clean causal identification is possible from this environmental family alone. This is not a failure of the framework but a diagnostic: it correctly identifies that the environmental suite is fundamentally insufficient to distinguish causal from confounded structure, and that additional structural assumptions (e.g., temporal precedence or instrumental variables) are needed.

\subsection{Connections to Optimal Transport}

The Wasserstein geometry on the space of probability measures provides an alternative Riemannian structure on $\cM$ where the tangent space has a natural interpretation in terms of transport plans. In the Wasserstein metric, the environmental Lie algebra would act via pushforward of transport maps, and the causal subspace would correspond to transport-invariant functionals. Characterizing $\Inv(\g)$ in Wasserstein geometry — where geodesics have a direct probabilistic interpretation as distributional interpolations — is a natural extension that may yield tighter identifiability conditions in settings where distributional robustness is the primary concern.

%% ─────────────────────────────────────────────────────────────────────────────
\section{Summary and Open Problems}
\label{sec:conclusion}
%% ─────────────────────────────────────────────────────────────────────────────

\subsection{The Chain of Equivalences}

The Invariant Manifold Transform establishes the following complete chain of equivalences for a smooth function $f$:
\begin{align}
f \text{ is causal} \;&\iff\; f \in \Inv(\g) \nonumber\\
&\iff\; f \text{ is constant on leaves of } \cF \nonumber\\
&\iff\; \nabla f \in \cH \nonumber\\
&\iff\; S_e(p) \nabla f = 0 \;\forall e \in \cE \;\text{(Theorem~\ref{thm:null_space_id}, under Spanning)} \nonumber\\
&\iff\; \Pi_{\ker(\hat S)} \nabla f = \nabla f \;\text{(computationally verifiable)}.
\label{eq:chain}
\end{align}
The final step bridges abstract geometry to an implementable algorithm. Table~\ref{tab:results} summarizes all main results.

\begin{table}[t]
\centering
\small
\begin{tabular}{llp{5cm}}
\toprule
\textbf{Result} & \textbf{Section} & \textbf{Content} \\
\midrule
Proposition~\ref{prop:hausdorff} & \S\ref{sec:foliation} & Properness $\Rightarrow$ Hausdorff quotient manifold \\
Lemma~\ref{lem:foliation} & \S\ref{sec:foliation} & Invariant quotient manifold exists; $\LP_{V_e}(\pi^* g_\pi) = 0$ \\
Theorem~\ref{thm:spectral_gap} & \S\ref{sec:lie} & Spectral Gap: dimension, identifiability, destruction threshold \\
Theorem~\ref{thm:bridge}(i) & \S\ref{sec:bridge} & Bridge (inclusion): $\cH_p \subseteq \bigcap_e \Null(S_e)$ always \\
Theorem~\ref{thm:bridge}(ii) & \S\ref{sec:bridge} & Bridge (equality): holds iff Spanning Condition satisfied \\
Theorem~\ref{thm:null_space_id} & \S\ref{sec:bridge} & Full null-space identification (complete, non-circular proof) \\
Proposition~6.4 & \S\ref{sec:bridge} & Spanning Condition is generic (measure-zero failure set) \\
Theorem~\ref{thm:stability} & \S\ref{sec:proxy} & Structural Stability: quantitative proxy exclusion bound \\
Theorem~\ref{thm:convergence} & \S\ref{sec:algorithm} & CRGD consistency and geometric convergence rate \\
\bottomrule
\end{tabular}
\caption{Summary of main results.}
\label{tab:results}
\end{table}

\subsection{Open Problems}

\paragraph{Finite-sample theory.} The convergence rate in Theorem~\ref{thm:convergence}(ii) has an additive bias term $C\sigma^2/N$ from finite-sample score estimation. A complete minimax analysis of the excess risk as a function of $N$, $|\cE|$, $d$, and $\delta_0$ is an important open problem.

\paragraph{Relaxing involutivity.} Assumption~\ref{ass:involutive} is sufficient but not necessary. An approximate involutivity condition $\|[X, Y] - Z\|_g \leq \varepsilon$ for some $Z \in \cD$ would yield an $\varepsilon$-approximate foliation and a corresponding perturbation of the causal subspace. Quantifying this perturbation in terms of $\varepsilon$ would extend the framework to settings where the environmental family only approximately generates a Lie algebra.

\paragraph{DAG structure and the Lie algebra.} When environments correspond to interventions on specific nodes of a DAG, the structure of $\g$ should reflect the graph (e.g., the Markov blanket of the target variable). A precise correspondence between the graphical criterion of $d$-separation and the algebraic criterion $f \in \Inv(\g)$ remains to be developed.

\paragraph{Non-parametric and infinite-dimensional extensions.} Extending from finite-dimensional $\theta \in \bR^d$ to infinite-dimensional function spaces (neural networks, reproducing kernel Hilbert spaces) requires replacing the finite-dimensional Fisher manifold with an infinite-dimensional statistical manifold~\cite{amari2000}, where the Frobenius theorem requires additional Banach-space regularity.

\paragraph{Wasserstein geometry.} Characterizing $\Inv(\g)$ in the Wasserstein metric, where the tangent space has a natural interpretation in terms of transport plans, may yield tighter identifiability conditions under distributional robustness objectives.

%% ─────────────────────────────────────────────────────────────────────────────
\section*{Conclusion}
%% ─────────────────────────────────────────────────────────────────────────────

The Invariant Manifold Transform reframes causal representation learning by identifying causality with algebraic invariance under a continuous symmetry group — in precise analogy with Noether's theorem in physics. The causal mechanism is not defined by graph structure, temporal precedence, or intervention semantics, but purely by its immunity to environmental perturbation:
\[
\boxed{\cC = \Inv(\g).}
\]
This principle, together with the complete chain of equivalences~\eqref{eq:chain}, transforms causal identification from a combinatorial search over graph space into a spectral computation on a sequence of Fisher information matrices. The mathematical contributions — a non-circular proof of null-space identification, the Spectral Gap Theorem with a sharp destruction threshold, the Structural Stability bound on proxy exclusion, and the Geodesic Projection Algorithm with linear convergence — together constitute a complete, rigorous, and computationally tractable theory of symmetry-based causal learning.

%% ─────────────────────────────────────────────────────────────────────────────
\begin{thebibliography}{99}
\bibitem{noether1918}
E.\ Noether, ``Invariante Variationsprobleme,'' \textit{Nachr.\ Ges.\ Wiss.\ Göttingen, Math.-Phys.\ Kl.}, pp.\ 235--257, 1918.

\bibitem{peters2016}
J.\ Peters, P.\ Bühlmann, and N.\ Meinshausen, ``Causal inference by using invariant prediction: identification and confidence intervals,'' \textit{J.\ Royal Statistical Society: Series B}, vol.\ 78, no.\ 5, pp.\ 947--1012, 2016.

\bibitem{arjovsky2019}
M.\ Arjovsky, L.\ Bottou, I.\ Gulrajani, and D.\ Lopez-Paz, ``Invariant Risk Minimization,'' \textit{arXiv:1907.02893}, 2019.

\bibitem{rosenfeld2020}
E.\ Rosenfeld, P.\ Ravikumar, and A.\ Risteski, ``The Risks of Invariant Risk Minimization,'' in \textit{Proc.\ ICLR 2021}.

\bibitem{amari2000}
S.\ Amari and H.\ Nagaoka, \textit{Methods of Information Geometry}, AMS/Oxford University Press, 2000.

\bibitem{lee2013}
J.\ M.\ Lee, \textit{Introduction to Smooth Manifolds}, 2nd ed., Springer, 2013.

\bibitem{palais1961}
R.\ S.\ Palais, ``On the existence of slices for actions of non-compact Lie groups,'' \textit{Ann.\ Math.}, vol.\ 73, no.\ 2, pp.\ 295--323, 1961.

\bibitem{kushner2003}
H.\ Kushner and G.\ G.\ Yin, \textit{Stochastic Approximation and Recursive Algorithms and Applications}, 2nd ed., Springer, 2003.

\bibitem{bott1982}
R.\ Bott and L.\ W.\ Tu, \textit{Differential Forms in Algebraic Topology}, Springer, 1982.
\end{thebibliography}

\end{document}